\chapter{Literature Review}

\section{Introduction to Skin Tone Detection}

Skin tone detection and classification have been active areas of research in computer vision, dermatology, and cosmetics for several decades. The ability to accurately detect and classify skin tones has applications ranging from medical diagnosis to personalized product recommendations. This chapter reviews the existing literature, technologies, and methodologies related to skin tone detection, color science, and cosmetic recommendation systems.

\section{Skin Tone Classification Systems}

\subsection{Fitzpatrick Skin Type Scale}

The Fitzpatrick scale, developed by dermatologist Thomas B. Fitzpatrick in 1975, is one of the most widely used skin classification systems in dermatology. It categorizes skin into six types (I-VI) based on response to UV radiation:

\begin{itemize}
    \item Type I: Very fair, always burns, never tans
    \item Type II: Fair, usually burns, tans minimally
    \item Type III: Medium, sometimes burns, tans uniformly
    \item Type IV: Olive, rarely burns, tans easily
    \item Type V: Brown, very rarely burns, tans very easily
    \item Type VI: Dark brown to black, never burns, tans very easily
\end{itemize}

While useful for dermatological purposes, the Fitzpatrick scale has significant limitations for cosmetic applications. It was developed primarily for Caucasian skin and does not adequately represent the full spectrum of human skin tones, particularly deeper tones. The scale's focus on UV response rather than visual appearance makes it less suitable for foundation matching.

\subsection{Monk Skin Tone Scale}

The Monk Skin Tone (MST) scale, developed by Dr. Ellis Monk at Harvard University, addresses the limitations of traditional classification systems. Published in 2019 and adopted by major technology companies including Google, the MST scale provides a more inclusive 10-class system:

\begin{table}[h]
\centering
\begin{tabular}{|c|l|l|}
\hline
\textbf{Class} & \textbf{Label} & \textbf{Description} \\
\hline
MST-01 & Porcelain & Very fair, cool/pink undertones \\
MST-02 & Fair & Fair, warm/peachy undertones \\
MST-03 & Light & Light beige \\
MST-04 & Light-Medium & Light medium \\
MST-05 & Medium & Medium beige \\
MST-06 & Medium-Dark & Warm medium \\
MST-07 & Tan & Tan/olive \\
MST-08 & Brown & Brown \\
MST-09 & Deep Brown & Deep brown \\
MST-10 & Deep & Deep/ebony \\
\hline
\end{tabular}
\caption{Monk Skin Tone Scale Classification}
\end{table}

The MST scale provides better representation for deeper skin tones, with four classes (MST-07 to MST-10) dedicated to brown and deep skin tones compared to only one or two in traditional systems. This makes it particularly suitable for cosmetic applications where accurate shade matching across the full spectrum is essential.

\subsection{Individual Typology Angle (ITA°)}

The Individual Typology Angle (ITA°) is a quantitative measure of skin tone based on CIELAB color space values. Introduced by Chardon et al. (1991), ITA° is calculated using the formula:

\begin{equation}
ITA° = \arctan\left(\frac{L^* - 50}{b^*}\right) \times \frac{180}{\pi}
\end{equation}

where $L^*$ is the lightness component and $b^*$ is the yellow-blue component in CIELAB color space. ITA° provides a continuous measure that can be mapped to discrete classification systems. Research has established threshold ranges for different skin tone categories, making it a valuable tool for automated classification.

\section{Color Science and Color Spaces}

\subsection{RGB Color Space}

RGB (Red, Green, Blue) is the most common color space for digital images, representing colors as combinations of red, green, and blue light. While intuitive and widely used, RGB is device-dependent and not perceptually uniform, meaning that equal numerical differences do not correspond to equal perceived color differences. This limitation makes RGB unsuitable for accurate color matching applications.

\subsection{CIELAB Color Space}

The CIELAB color space, developed by the Commission Internationale de l'Éclairage (CIE) in 1976, is designed to be perceptually uniform and device-independent. It consists of three components:

\begin{itemize}
    \item $L^*$: Lightness (0 = black, 100 = white)
    \item $a^*$: Red-green axis (negative = green, positive = red)
    \item $b^*$: Yellow-blue axis (negative = blue, positive = yellow)
\end{itemize}

CIELAB is particularly valuable for skin tone analysis because:
\begin{enumerate}
    \item It separates lightness from chromaticity, allowing independent analysis of tone and undertone
    \item It is perceptually uniform, meaning numerical differences correspond to perceived differences
    \item It is device-independent, providing consistent results across different cameras and displays
    \item It forms the basis for industry-standard color difference metrics
\end{enumerate}

\subsection{HSV and YCbCr Color Spaces}

HSV (Hue, Saturation, Value) and YCbCr (Luma, Blue-difference, Red-difference) color spaces are commonly used for skin detection in computer vision applications. Research has shown that skin pixels cluster in specific regions of these color spaces, enabling rule-based segmentation:

\begin{itemize}
    \item HSV: Skin typically has hue values between 0-50° and saturation between 0.2-0.6
    \item YCbCr: Skin pixels satisfy conditions like $77 \leq Cb \leq 127$ and $133 \leq Cr \leq 173$
\end{itemize}

These color spaces are particularly effective for skin segmentation because they separate chromaticity from luminance, making them robust to lighting variations.

\subsection{Color Difference Metrics}

\subsubsection{CIEDE2000}

CIEDE2000 is the current industry standard for measuring color difference, published by CIE in 2000. It improves upon earlier metrics (ΔE*ab, ΔE*94) by incorporating:

\begin{itemize}
    \item Weighting functions for lightness, chroma, and hue
    \item Rotation term for blue region
    \item Compensation for neutral colors
    \item Perceptual uniformity corrections
\end{itemize}

The CIEDE2000 formula is complex but provides the most accurate measure of perceived color difference. A ΔE value below 1.0 is considered imperceptible, below 2.0 is excellent, and below 3.0 is acceptable for most applications. Research in cosmetics typically targets ΔE < 3.0 for foundation matching.

\section{Computer Vision Techniques}

\subsection{Face Detection}

Face detection is a critical first step in skin tone analysis. Several approaches have been developed:

\subsubsection{Haar Cascade Classifiers}

Introduced by Viola and Jones (2001), Haar Cascade classifiers use machine learning to detect faces based on Haar-like features. The method is:
\begin{itemize}
    \item Fast and efficient (real-time performance on CPU)
    \item Robust to variations in pose and lighting
    \item Available in OpenCV with pre-trained models
    \item Suitable for frontal face detection
\end{itemize}

Limitations include reduced accuracy for non-frontal faces and sensitivity to occlusions.

\subsubsection{Multi-task Cascaded Convolutional Networks (MTCNN)}

MTCNN, proposed by Zhang et al. (2016), uses deep learning for joint face detection and alignment. It provides:
\begin{itemize}
    \item Higher accuracy than Haar Cascades
    \item Facial landmark detection
    \item Better handling of various poses and expressions
    \item Confidence scores for detections
\end{itemize}

However, MTCNN requires more computational resources and may need GPU acceleration for real-time performance.

\subsection{Skin Segmentation}

Skin segmentation aims to identify skin pixels in an image, separating them from background, hair, eyes, and other non-skin regions.

\subsubsection{Rule-Based Methods}

Rule-based methods use color space thresholds to identify skin pixels. Common approaches include:

\begin{itemize}
    \item HSV thresholding: $0° \leq H \leq 50°$, $0.2 \leq S \leq 0.6$
    \item YCbCr thresholding: $77 \leq Cb \leq 127$, $133 \leq Cr \leq 173$
    \item Combined methods using multiple color spaces
\end{itemize}

Advantages include simplicity, speed, and no training requirement. Limitations include sensitivity to lighting and difficulty handling diverse skin tones.

\subsubsection{Machine Learning Methods}

Machine learning approaches train classifiers on labeled skin/non-skin pixels:

\begin{itemize}
    \item Gaussian Mixture Models (GMM)
    \item Support Vector Machines (SVM)
    \item Random Forests
    \item Neural Networks
\end{itemize}

These methods can achieve higher accuracy but require training data and more computational resources.

\subsubsection{Deep Learning Methods}

Deep learning approaches, particularly U-Net and its variants, have achieved state-of-the-art results in skin segmentation:

\begin{itemize}
    \item U-Net: Encoder-decoder architecture with skip connections
    \item SegNet: Similar architecture with pooling indices
    \item DeepLab: Uses atrous convolution for multi-scale features
\end{itemize}

Deep learning methods provide the highest accuracy but require large annotated datasets, GPU hardware, and longer training times.

\section{Existing Skin Tone Detection Systems}

\subsection{Academic Research}

Several academic papers have addressed skin tone detection and classification:

\subsubsection{Deep Learning Approaches}

Recent research has focused on deep learning methods:

\begin{itemize}
    \item \textbf{EfficientNet-based Classification:} Researchers have used EfficientNet architectures for MST classification, achieving 92-94\% accuracy on test sets. However, these methods require thousands of annotated images and GPU training.
    
    \item \textbf{CNN + SVM Hybrid:} Some studies combine CNN feature extraction with SVM classification, reporting ΔE values around 3.0-3.5 for foundation matching.
    
    \item \textbf{ResNet-based Systems:} ResNet architectures have been applied to skin tone classification with reported accuracies of 88-90\% on Fitzpatrick scale classification.
\end{itemize}

Common limitations include:
\begin{itemize}
    \item High computational requirements
    \item Need for large annotated datasets
    \item Limited foundation databases (typically 50-80 shades)
    \item Lack of production-ready implementations
\end{itemize}

\subsubsection{Traditional Computer Vision Approaches}

Classical computer vision methods have also been explored:

\begin{itemize}
    \item ITA°-based classification achieving 85-90\% accuracy
    \item Rule-based segmentation with morphological operations
    \item Color histogram analysis for undertone detection
\end{itemize}

These methods are faster and require no training but may have lower accuracy for challenging cases.

\subsection{Commercial Systems}

Several companies have developed skin tone detection systems:

\subsubsection{L'Oréal Shade Finder}

L'Oréal's mobile app uses AI to analyze selfies and recommend foundation shades. Features include:
\begin{itemize}
    \item Real-time camera analysis
    \item Integration with product catalog
    \item Virtual try-on capabilities
\end{itemize}

Technical details are proprietary, but the system likely uses deep learning models trained on extensive datasets.

\subsubsection{Sephora Virtual Artist}

Sephora's app provides virtual makeup try-on and shade matching:
\begin{itemize}
    \item Augmented reality overlay
    \item Multiple product categories
    \item Integration with e-commerce
\end{itemize}

The system focuses on user experience and visual appeal rather than color accuracy.

\subsubsection{Perfect Corp. YouCam Makeup}

YouCam Makeup offers comprehensive virtual beauty features:
\begin{itemize}
    \item Real-time face tracking
    \item Multiple makeup products
    \item Skin analysis features
\end{itemize}

The system uses proprietary AI models and extensive product databases.

\subsection{Limitations of Existing Systems}

Current systems, both academic and commercial, have several limitations:

\begin{enumerate}
    \item \textbf{Accuracy:} Many systems achieve ΔE values of 3.0-5.0, which is acceptable but not excellent. Deeper skin tones often have higher errors.
    
    \item \textbf{Accessibility:} Deep learning systems require expensive GPU hardware and are not accessible to individual developers or small businesses.
    
    \item \textbf{Training Data:} Most systems require thousands of annotated images, which are expensive and time-consuming to collect.
    
    \item \textbf{Database Coverage:} Foundation databases are often limited to 50-100 shades, with inadequate representation of deeper tones.
    
    \item \textbf{Production Readiness:} Academic prototypes often lack user-friendly interfaces, documentation, and deployment options.
    
    \item \textbf{Inclusivity:} Many systems use outdated classification scales (Fitzpatrick) that do not adequately represent diverse skin tones.
\end{enumerate}

\section{Research Gap and Opportunity}

Based on the literature review, several gaps and opportunities have been identified:

\subsection{Gaps in Current Research}

\begin{enumerate}
    \item \textbf{Accessibility Gap:} Most high-accuracy systems require GPU hardware and extensive training, making them inaccessible to many users.
    
    \item \textbf{Database Gap:} Existing foundation databases are limited in size and diversity, particularly for deeper skin tones.
    
    \item \textbf{Implementation Gap:} Few research papers provide production-ready implementations with user-friendly interfaces.
    
    \item \textbf{Inclusivity Gap:} Many systems still use outdated classification scales that do not adequately represent diverse populations.
    
    \item \textbf{Documentation Gap:} Academic papers often lack comprehensive documentation and reproducible code.
\end{enumerate}

\subsection{Opportunities for Innovation}

\begin{enumerate}
    \item \textbf{Classical Methods:} Combining well-established computer vision and color science techniques can achieve excellent results without requiring training data or GPU hardware.
    
    \item \textbf{Comprehensive Database:} Building a large, diverse foundation database with accurate CIELAB values can significantly improve matching accuracy.
    
    \item \textbf{Inclusive Classification:} Adopting the Monk Skin Tone scale provides better representation and aligns with industry trends toward inclusivity.
    
    \item \textbf{Production-Ready System:} Developing a complete system with multiple interfaces (CLI, Web, API) and comprehensive documentation fills a significant gap.
    
    \item \textbf{Open Source:} Making the system open source enables community contributions, improvements, and wider adoption.
\end{enumerate}

\section{Theoretical Foundation}

\subsection{Color Perception and Matching}

Human color perception is complex and influenced by multiple factors:

\begin{itemize}
    \item \textbf{Lighting:} The same object appears different under different illuminants (metamerism)
    \item \textbf{Surrounding Colors:} Adjacent colors affect perception (simultaneous contrast)
    \item \textbf{Individual Variation:} Color perception varies between individuals
    \item \textbf{Context:} Expectations and context influence color judgment
\end{itemize}

For foundation matching, the goal is to minimize the perceived color difference between the foundation and skin under typical viewing conditions (daylight, indoor lighting). CIEDE2000 provides the best mathematical model of perceived color difference, making it the appropriate metric for this application.

\subsection{White Balance and Color Constancy}

Digital cameras capture images under various lighting conditions, which can significantly affect color appearance. White balance correction aims to remove color casts and make colors appear as they would under neutral illumination. The Gray World algorithm, used in ShadeAI, assumes that the average color in an image should be neutral gray. While simple, it is effective for most photographs and requires no training or calibration.

\subsection{Skin Tone and Undertone}

Skin tone refers to the overall lightness or darkness of skin, while undertone refers to the underlying hue:

\begin{itemize}
    \item \textbf{Warm Undertone:} Yellow, golden, or peachy hues (high $b^*$ in CIELAB)
    \item \textbf{Cool Undertone:} Pink, red, or bluish hues (low $b^*$, high $a^*$)
    \item \textbf{Neutral Undertone:} Balanced mix of warm and cool
\end{itemize}

Matching both tone and undertone is essential for natural-looking foundation. A foundation that matches in lightness but has the wrong undertone will appear unnatural.

\section{Summary}

This literature review has examined the current state of research and technology in skin tone detection, color science, and cosmetic recommendation systems. Key findings include:

\begin{enumerate}
    \item The Monk Skin Tone scale provides a more inclusive classification system than traditional scales
    \item CIELAB color space and CIEDE2000 metrics are essential for accurate color matching
    \item Both classical computer vision and deep learning approaches have been explored, each with advantages and limitations
    \item Existing systems have gaps in accessibility, database coverage, and production readiness
    \item There is an opportunity to develop a practical system that achieves research-grade accuracy using classical methods
\end{enumerate}

The next chapter describes the system design and methodology for ShadeAI, which addresses these gaps and leverages the theoretical foundations established in this review.
